\begin{abstract}
Robots ask people questions to resolve ambiguity, to elicit preferences and to check trust, and they treat the answers as readings of fixed beliefs. Human answers, however, depend on question order, and asking can change what a person thinks next. This paper builds an order-aware human model for robot questioning and tests it against classical alternatives. A three-dimensional quantum-like (quantum-probability) model with projector ranks is compared with an order-free Bayesian model and two anchoring models, and an open-system trust model is compared with a Markov model, in simulated populations for three robot domains: object clarification, trust and hand-over, and preference elicitation. On published data, the quantum-like model fits the Clinton--Gore order effect as well as anchoring does, whereas the Bayesian model cannot. In simulation, the quantum-like model is identified from 400 people per question order in 93\% of homogeneous data sets and is then the best predictor of answers in the unasked order, but individual differences reduce its advantage, and when people follow a classical order effect it is the worst predictor, about ten times worse than the Bayesian model. For recovering unprimed answers, randomizing question order beats model-based correction under every generator. Under open-system dynamics, a midway trust question lowers final trust by 0.076; detecting this needs about 1,600 people per condition, and a robot that always asks cannot learn it. Both models run as quantum circuits; under the noise model of a current IBM processor, answer distributions deviate by about 0.04 in total variation and the QQ equality holds to within 0.012. Quantum-like models are therefore a useful but conditional part of a robot's model of people: they should be weighted against classical models by evidence, not used alone.
\end{abstract}

\begin{IEEEkeywords}
Human--robot interaction, question order effects, quantum cognition, quantum computing, trust, user modeling.
\end{IEEEkeywords}
